\par\medskip
\noindent
2.3. Supongamos de nuevo que $S$ es henseliano y que
$\ell\ne p=\operatorname{char}k$. Utilizamos ahora la inclusión
\begin{equation}\tag{2.3.1}
 T_\ell(T_o)\hookrightarrow T_\ell(A_o^o)
\end{equation}
de (2.1.6), junto con la equivalencia de categorías recordada en 2.2.4.
Obtenemos un subhaz $\ell$-ádico constante torcido canónico de
$T_\ell(A^o)^f$, llamado la \underline{parte tórica de}
$T_\ell(A^o)$,
\begin{equation}\tag{2.3.2}
 T_\ell(A^o)^t\hookrightarrow T_\ell(A^o)^f
 \hookrightarrow T_\ell(A^o),
\end{equation}
caracterizado por la condición de que su fibra especial sea la
inclusión (2.3.1):
\begin{equation}\tag{2.3.3}
 T_\ell(A^o)^t\times_Ss=T_\ell(T_o)
 \hookrightarrow T_\ell(A_o^o).
\end{equation}
Más adelante daremos una interpretación de la parte tórica de
$T_\ell(A^o)$ como el módulo de Tate $T_\ell(T)$ de un toro $T$ sobre
$S$ que levanta el toro $T_o$ sobre $k$. Esta interpretación permitirá
extender la definición de la parte tórica al caso $\ell=p$. Por el
momento, y para fijar ideas, nos limitamos al caso $\ell\ne p$. En este
caso, la parte tórica de $T_\ell(A^o)$ puede describirse mediante el
$\mathbb Z_\ell$-submódulo correspondiente de
$T_\ell(A(\overline K))$:
\begin{equation}\tag{2.3.4}
 T_\ell(A(\overline K))^t\hookrightarrow
 T_\ell(A(\overline K))^f\hookrightarrow
 T_\ell(A(\overline K)).
\end{equation}
Es estable bajo la acción de $\pi$; su primer término se define como
$T_\ell(A^o)^t(\overline K)$ y, como cabría esperar, se denomina de
nuevo el \underline{submódulo tórico} del módulo de Tate
$T_\ell(A(\overline K))$. Así, este módulo de Tate lleva la filtración
canónica de tres escalones (2.3.4). Aplicando las mismas consideraciones
al esquema abeliano dual $A'_K$, obtenemos la filtración análoga de tres
escalones
\begin{equation}\tag{2.3.5}
 T_\ell(A'(\overline K))^t\hookrightarrow
 T_\ell(A'(\overline K))^f\hookrightarrow
 T_\ell(A'(\overline K)).
\end{equation}
Los términos intermedios de (2.3.4) y (2.3.5), las «partes fijas», se
describen en términos de las estructuras de $\pi$-módulos de los
módulos de Tate de $A(\overline K)$ y $A'(\overline K)$ como sus
submódulos de invariantes bajo el grupo de inercia $I$
(\sisi{1.6.5}{2.2.5}). Las partes tóricas quedan entonces completamente
descritas en términos de las partes fijas y del emparejamiento
$\varphi$ de (1.0.2), mediante el teorema siguiente (del que ya aparece
un caso particular en [11]).

\medskip
\noindent
\underline{Teorema} 2.4. (\underline{Teorema de ortogonalidad}.)
\underline{Sea $S$ un rasgo henseliano, sea $A_K$ un esquema abeliano
sobre su cuerpo de fracciones $K$, sea $A'_K$ el esquema abeliano dual,
y sea $\ell$ un primo distinto de la característica residual de $S$.
Los módulos de Tate $T_\ell(A(\overline K))$ y
$T_\ell(A'(\overline K))$ llevan entonces, respectivamente, las
filtraciones canónicas (2.3.4) y (2.3.5) por la «parte fija» y la «parte
tórica». La parte tórica $T_\ell(A(\overline K))^t$ es la intersección
de la parte fija $T_\ell(A(\overline K))^f$ con el ortogonal de la parte
fija $T_\ell(A'(\overline K))^f=T_\ell(A'(\overline K))^I$ respecto del
emparejamiento canónico $\varphi$ de (1.0.3). Simétricamente, la parte
tórica $T_\ell(A'(\overline K))^t$ es la intersección de
$T_\ell(A'(\overline K))^f$ con el ortogonal de la parte fija
$T_\ell(A(\overline K))^f=T_\ell(A(\overline K))^I$. (Aquí el
superíndice $I$ denota el submódulo de invariantes bajo la acción del
grupo de inercia $I$.)}

\smallskip
\noindent
\underline{Demostración.}\quad a) Demostremos que
$T_\ell(A(\overline K))^t$ es ortogonal a
$T_\ell(A'(\overline K))^f$. Para ello, interpretamos la restricción
de $\varphi$ a las partes fijas de los módulos de Tate mediante la
biextensión canónica $W^o=W$ de (1.4.4). Esta es el prolongamiento
canónico de la biextensión $W_K$ de $(A_K\sisi{;}{,}A'_K)$ por $\mathbb G_{mK}$
a una biextensión de $(A^o,A'{}^o)$ por $\mathbb G_{mS}$. Según
VIII~2.2, $W$ da lugar a un emparejamiento
\begin{equation}\tag{2.4.1}
 \varphi_{\sisi{w}{W}}:T_\ell(A^o)\times T_\ell(A'{}^o)
 \longrightarrow T_\ell(\mathbb G_{mS}),
\end{equation}
del que se obtiene el emparejamiento $\varphi$ de (1.0.3) pasando a los
puntos con valores en $\overline K$. Se sigue que la restricción de
$\varphi$ a las partes fijas consideradas en 2.4 se obtiene, pasando a
los puntos con valores en $\overline K$, del emparejamiento
\begin{equation}\tag{2.4.2}
 T_\ell(A^o)^f\times T_\ell(A'{}^o)^f
 \longrightarrow T_\ell(\mathbb G_{mS}),
\end{equation}
restricción de (2.4.1). Puesto que los dos factores del miembro
izquierdo de (2.4.2) son ahora haces $\ell$-ádicos constantes torcidos,
este emparejamiento puede estudiarse sobre las fibras especiales. Las
fibras de ambos factores son precisamente los módulos de Tate
$T_\ell(A_o^o)$ y $T_\ell(A'{}_o^o)$ (2.2.3.3), y el emparejamiento
inducido sobre las fibras especiales por (2.4.2) es
\begin{equation}\tag{2.4.3}
 T_\ell(A_o^o)\times T_\ell(A'{}_o^o)
 \longrightarrow T_\ell(\mathbb G_{mk}),
\end{equation}
obtenido de la biextensión $W_o$ de $(A_o^o,A'{}_o^o)$ por
$\mathbb G_{mk}$, restricción de $W$ a las fibras especiales.

Recordemos de 2.3 que la parte tórica de $T_\ell(A(\overline K))$ se
obtiene pasando a los puntos con valores en $\overline K$ en el subhaz
$\ell$-ádico constante torcido $T_\ell(A^o)^t$ de $T_\ell(A^o)^f$.
Así, la relación de ortogonalidad que se quiere demostrar expresa que
el emparejamiento inducido por (2.4.2),
\begin{equation}\tag{2.4.4}
 T_\ell(A^o)^t\times T_\ell(A'{}^o)^f
 \longrightarrow T_\ell(\mathbb G_{mS}),
\end{equation}
es idénticamente nulo. Basta comprobarlo para el emparejamiento
correspondiente sobre las fibras especiales, que por la definición de
la parte tórica es
\begin{equation}\tag{2.4.5}
 T_\ell(T_o)\times T_\ell(\sisi{A_o^{\,o}}{A_o^{\prime\,o}})
 \longrightarrow T_\ell(\mathbb G_{mk}),
\end{equation}
inducido por (2.4.3). Según VIII~4.10, este emparejamiento es
necesariamente nulo, lo que demuestra la ortogonalidad afirmada.

\smallskip
\noindent
b) Demostremos que la inclusión
\begin{equation}\tag{2.4.6}
 T_\ell(A(\overline K))^t\subset
 T_\ell(A(\overline K))^f\cap
 \bigl(T_\ell(A'(\overline K))^f\bigr)^\perp
\end{equation}
es una igualdad, donde $\perp$ denota el ortogonal. Equivalentemente,
en el emparejamiento (2.4.2), el anulador del segundo factor ---que,
como acabamos de ver, contiene la parte tórica $T_\ell(A^o)^t$--- es
igual a esta parte tórica. Puesto que se trata de nuevo de relaciones
entre haces $\ell$-ádicos constantes torcidos, basta comprobar la
afirmación sobre las fibras especiales: a saber, que la inclusión ya
obtenida \sisi{(1.10)}{(2.4.7)},
\begin{equation}\tag{2.4.7}
 T_\ell(T_o)\subset T_\ell(A'{}_o^o)^\perp
\end{equation}
es una igualdad. En otras
palabras, debemos comprobar que la biextensión $W_o$ de
$(A_o^o,A'{}_o^o)$ es \underline{no degenerada por la izquierda}
(VIII~4.11). Pero esto es precisamente lo que ya observamos en 1.5,
como consecuencia de un resultado de la tesis de Raynaud.
